\subsection{EXPLICIT FORM OF INNER PRODUCTS IN THE CASE OF A FINITELY GENERATED FREE MODULE}

Let \(\mathbf{M}\) be a finitely generated free A-module, \((e_i)_{1\leq i\leq n}\) a basis of \(\mathbf{M}\) and \((e_i^*)_{1\leq i\leq n}\) the dual basis of \(\mathbf{M}^*\) . For every finite sequence \(s = (i_1,\ldots ,i_p)\) of elements of \([1,n]\) , let \(e_s = e_{i_1}\otimes e_{i_2}\otimes \dots \otimes e_{i_p}\) (resp. \(e_s^* = e_{i_1}^*\otimes \dots \otimes e_{i_p}^*\) ). We know (§ 5, no. 5, Theorem 1) that the \(e_s\) form a basis of the A-module \(\mathsf{T}(\mathbf{M})\) and the \(e_s^*\) a basis of the A-module \(\mathsf{T}(\mathbf{M}^*)\) . If \(s,t\) are two finite sequences of elements of \([1,n]\) , let \(s.t\) denote the sequence obtained as follows: if \(s = (i_1,\ldots ,i_p)\) and \(t = (j_1,\ldots ,j_q)\) , \(s.t\) is the sequence \((i_1,\ldots ,i_p,j_1,\ldots ,j_q)\) with \(p + q\) terms. Then \(e_{s,t} = e_s\otimes e_t\) . It then follows from (66) that

\[
\left\{ \begin{array}{l} e _ {s} \llcorner e _ {t} ^ {*} = 0 \quad \text { if   } s \text {   is   not   of   the   form   } t. u \\ e _ {t. u} \llcorner e _ {t} ^ {*} = e _ {u}. \end{array} \right.\tag{72}
\]

Similarly, the symmetric algebra \(\mathbf{S}(\mathbf{M})\) has as basis the set of monomials \(e^{\alpha}\) with \(\alpha \in \mathbf{N}^{n}\) (§ 6, no. 6, Theorem 1) and \(\mathbf{S}(\mathbf{M}^{*})\) the set of monomials \(e^{*\alpha}\) with \(\alpha \in \mathbf{N}^{n}\) ; recall (no. 5) that \(u_{\alpha}\) , for \(|\alpha| = k\) , denotes the element of the basis of \((\mathbf{S}^{k}(\mathbf{M}))^{*}\) , dual to the basis \((e^{\alpha})_{|\alpha| = k}\) of \(\mathbf{S}^{k}(\mathbf{M})\) ; the \(u_{\alpha}\) , for \(\alpha \in \mathbf{N}^{n}\) , therefore form a basis of \(\mathbf{S}(\mathbf{M})^{*\mathrm{gr}}\) . The definition of right inner product by \(e^{\beta}\) in \(\mathbf{S}(\mathbf{M})^{*\mathrm{gr}}\) as the transpose of multiplication by \(e^{\beta}\) in \(\mathbf{S}(\mathbf{M})\) then shows that

\[
\left\{ \begin{array}{l l} u _ {\alpha} \llcorner e ^ {\beta} = 0 & \text { if } \quad \alpha \not \geqslant \beta \\ u _ {\alpha} \llcorner e ^ {\beta} = u _ {\alpha - \beta} & \text { if } \quad \alpha \geqslant \beta . \end{array} \right.\tag{73}
\]

Similarly, since \(\mathbf{S}(\mathbf{M})\) is here canonically identified with the graded dual of \(\mathbf{S}(\mathbf{M})^{*\varepsilon r}\) , \(i(u_{\beta})\) is the graded transpose of multiplication by \(u^{\beta}\) in \(\mathbf{S}(\mathbf{M})^{*\varepsilon r}\) and hence from the multiplication table (33) (no. 5) of the basis \((u_{\alpha})\) we deduce that

\[
\left\{ \begin{array}{l l} e ^ {\alpha} \llcorner u _ {\beta} = 0 & \text { if } \quad \alpha \not \geqslant \beta \\ e ^ {\alpha} \llcorner u _ {\beta} = (\beta , \alpha - \beta) e ^ {\alpha - \beta} & \text { if } \quad \alpha \geqslant \beta . \end{array} \right.\tag{74}
\]

As for the right inner product of an element of \(\mathsf{S}(\mathbf{M})\) by an element of \(\mathsf{S}(\mathbf{M}^*)\) , the definition of this product (no. 9) and formula (34) of no. 5 allow us to deduce from (74) the formulae

\[
\left\{ \begin{array}{l l} e ^ {\alpha} \llcorner e ^ {* \beta} = 0 & \text { if } \quad \alpha \not \geqslant \beta \\ e ^ {\alpha} \llcorner e ^ {* \beta} = \frac {\alpha !}{(\alpha - \beta) !} e ^ {\alpha - \beta} & \text { if } \quad \alpha \geqslant \beta . \end{array} \right.\tag{75}
\]

There are analogous formulae for the inner product of an element of \(\mathsf{S}(\mathbf{M}^{*})\) by an element of \(\mathsf{S}(\mathbf{M})\) interchanging the roles of M and \(M^{*}\) (since \(M^{**}\) is here identified with M).

Remark. Being given the basis \((e_i)_{1\leq i\leq n}\) allows us to identify the algebra \(\mathsf{S}(\mathbf{M})\) with the polynomial algebra \(\mathbf{A}[\mathbf{X}_1,\ldots ,\mathbf{X}_n]\) (§ 6, no. 6); formula (75) shows that the inner product by \(e^{*\alpha}\) is just the differential operator \(\mathbf{D}^{\alpha} = \mathbf{D}_{1}^{\alpha_{1}}\mathbf{D}_{2}^{\alpha_{2}}\ldots \mathbf{D}_{n}^{\alpha_{n}}\) , where \(\mathbf{D}_i = \partial /\partial \mathbf{X}_i\) for \(1\leqslant i\leqslant n\) (§ 10, no. 11, Example).

Consider finally the exterior algebra \(\bigwedge (\mathbf{M})\) , which has as basis the set of elements \(e_{\mathbf{J}}\) , where \(\mathbf{J}\) runs through the set of subsets of the interval \([1, n]\) of \(\mathbf{N}\) (§ 7, no. 8, Theorem 1); similarly \(\bigwedge (\mathbf{M}^*)\) has as basis the elements \(e_{\mathbf{J}}^{*}\) . It follows from formula (68) of no. 9 that

\[
\left\{ \begin{array}{l l} e _ {\mathrm{K}} ^ {*} \lrcorner e _ {\mathrm{J}} = 0 & \text { if } \quad \mathrm{K} \nsubseteq \mathrm{J} \\ e _ {\mathrm{K}} ^ {*} \lrcorner e _ {\mathrm{J}} = (- 1) ^ {p (p - 1) / 2} \rho_ {\mathrm{K}, \mathrm{J} - \mathrm{K}} e _ {\mathrm{J} - \mathrm{K}} & \text { if } \quad \mathrm{K} \subset \mathrm{J} \quad \text { and } \quad p = \operatorname{Card} (\mathrm{K}), \end{array} \right.\tag{76}
\]

where \(\rho_{\mathbf{K},J - \mathbf{K}}\) is the number defined by formula (19) of §7, no. 8. There are analogous formulae with the roles of \(\mathbf{M}\) and \(\mathbf{M}^*\) interchanged.

